Inspectors are C++ classes which are usually bound to graph
labels. They are normally used to display the contents of a label in a
window on the screen.

E.g. the window label {\bf FFTclWindow} uses the class {\bf
FFTclGraphInsp} to display the contents of graph, node and egde
labels.

\begin{figure}[htb]
\noindent
\fbox{\epsfig{figure=inspector.eps,clip=}}
\caption{FFTclGraphInsp}
\end{figure}

\subsection{Using inspectors}

Inspectors are subclasses of {\bf FFInspector}. 
Inspectors are usually invoked by the graph label they are bound to
and stay on the screen until the label is deleted or the window is
dismissed by the user. See the corresponding graph label methods for
details on how to invoke inspectors.

\subsection{Writing inspectors}
\label{writing_inspectors}

If you write your own graph label, you might also want to write an
inspector for it to display its contents. To write an inspector you
have to subclass {\bf FFInspector} and implement the following
methods:

\begin{itemize}

\item{\bf FFInspector(FFGraphLabel *g)}

Initialize the inspector according to {\bf g}. {\bf g} is the label
the inspector is bound to.

\item{\bf virtual void show()}

Should display the inspector on the screen.

\item{\bf virtual void hide()}

Should remove the inspector from the screen.

\item{\bf virtual void inspect() = NULL}

Display the contents of the label the inspector is bound to, or
whatever you want to display.

\end{itemize}

If you want to use the tcl/tk widget set to display the inspector, you
might want to subclass {\bf FFTclInspector} which is a subclass of
{\bf FFInspector}.

